\documentclass[12pt,a4paper]{article}
\usepackage[utf8]{inputenc}
\usepackage[T1]{fontenc}
\usepackage[polish]{babel}
\usepackage{amsmath, amssymb, amsfonts}
\usepackage{geometry}
\geometry{a4paper, margin=2.5cm}

\begin{document}

\begin{center}
    \Large \textbf{Kolokwium z Analizy Matematycznej 1.2*} \\
    \vspace{0.2cm}
    \normalsize \textbf{25 kwietnia 2023, 16.00--19.00}
\end{center}

\vspace{0.8cm}

\noindent\textbf{Zadanie 1.} Wykaż, że dla dowolnych liczb dodatnich $a$, $b$, $c>0$ zachodzi nierówność:
\[
(a+b)^{c}\cdot(b+c)^{a}\cdot(a+c)^{b}\le\left(\frac{2}{3}(a+b+c)\right)^{(a+b+c)}
\]

\vspace{0.5cm}

\noindent\textbf{Zadanie 2.} Rozważmy funkcję $f:\mathbb{R}\rightarrow(0,+\infty)$ klasy $C^{1}$. Udowodnij, że istnieje $\xi\in[0,1]$ takie, że:
\[
e^{f^{\prime}(\xi)}f(0)^{f(\xi)}=f(1)^{f(\xi)}.
\]

\vspace{0.5cm}

\noindent\textbf{Zadanie 3.} Rozważmy funkcję różniczkowalną $f:\mathbb{R}\rightarrow\mathbb{R}$. Wykaż, że na dowolnym odcinku $[a,b]\subset\mathbb{R}$ funkcje $f^{\prime}$ oraz $|f^{\prime}|$ mają takie same wahania.

\vspace{0.5cm}

\noindent\textbf{Zadanie 4.} Rozważmy liczby $0<a<b$ i $p=\frac{1}{3}$. Udowodnij, że:
\[
\sqrt{ab}<\frac{b-a}{\ln b-\ln a}<\left(\frac{a^{p}+b^{p}}{2}\right)^{1/p}
\]
Czy stałą $p$ można zastąpić liczbą mniejszą?

\vspace{0.5cm}

\noindent\textbf{Zadanie 5.} Rozstrzygnij, czy istnieje funkcja ciągła $f:\mathbb{R}\rightarrow\mathbb{R}$ taka, że dla każdego $x\in\mathbb{R}$ istnieje granica (niewłaściwa) $f^{\prime}(x):=\lim_{h\rightarrow0}\frac{f(x+h)-f(x)}{h}$ przy czym $f^{\prime}(x)=+\infty$ lub $f^{\prime}(x)=-\infty$.

\end{document}